% Fuente: Auxiliar 3, sección «Método gráfico y soluciones básicas factibles».
\textbf{(a) Modelado como LP y Forma Estándar}
\vspace{2mm}

Definamos:
\[
x=\text{hectáreas de trigo},\quad y=\text{hectáreas de maíz}.
\]
El LP se plantea como:
\[
\begin{array}{rl}
\max & 200x+300y\\[1mm]
\text{sujeto a} & x+y \le 45,\\[1mm]
& 3x+2y \le 100,\\[1mm]
& 2x+4y \le 120,\\[1mm]
& x,y\ge 0.
\end{array}
\]
Para obtener la forma estándar se introducen variables de holgura \( s_1, s_2, s_3 \):
\[
\begin{array}{rl}
\max & 200x+300y\\[1mm]
\text{sujeto a} & x+y+s_1 = 45,\\[1mm]
& 3x+2y+s_2 = 100,\\[1mm]
& 2x+4y+s_3 = 120,\\[1mm]
& x,y,s_1,s_2,s_3 \ge 0.
\end{array}
\]

\textbf{(b) Resolución Gráfica y Soluciones Básicas Factibles}
\vspace{2mm}

La región factible está definida por las siguientes restricciones:
\begin{itemize}
    \item \( x+y\le 45 \),
    \item \( 3x+2y\le 100 \),
    \item \( 2x+4y\le 120 \),
    \item \( x,y\ge 0 \).
\end{itemize}

\begin{figure}
    \centering
    \includegraphics[width=0.7\linewidth]{ejercicios/Capitulo-3/ejercicio_3_3_region_factible.png}
    \caption{Región Factible}
    \label{ex:3:3:region-factible}
\end{figure}


Al resolver gráficamente, los vértices de la región factible se obtienen a partir de las intersecciones entre las restricciones activas. Los puntos extremos factibles son:
\begin{itemize}
    \item \((0,0)\): corresponde al origen, donde no se planta trigo ni maíz.

    \item \((0,30)\): al fijar \(x=0\), se tiene
    \[
    2(0)+4y\le 120 \quad \Rightarrow \quad y\le 30.
    \]
    Por lo tanto, el punto extremo sobre el eje \(y\) es \((0,30)\).

    \item \((20,20)\): intersección de
    \[
    3x+2y=100
    \]
    y
    \[
    2x+4y=120.
    \]
    Simplificando la segunda ecuación:
    \[
    x+2y=60.
    \]
    Luego, resolviendo el sistema:
    \[
    \begin{cases}
    3x+2y=100,\\
    x+2y=60,
    \end{cases}
    \]
    se obtiene:
    \[
    2x=40 \quad \Rightarrow \quad x=20.
    \]
    Sustituyendo en \(x+2y=60\):
    \[
    20+2y=60 \quad \Rightarrow \quad y=20.
    \]
    Por lo tanto, el punto es \((20,20)\).

    \item \(\left(\frac{100}{3},0\right)\): al fijar \(y=0\), se tiene
    \[
    3x+2(0)\le 100 \quad \Rightarrow \quad x\le \frac{100}{3}.
    \]
    Por lo tanto, el punto extremo sobre el eje \(x\) es
    \[
    \left(\frac{100}{3},0\right).
    \]
\end{itemize}

Es importante notar que las intersecciones \((10,35)\) y \((30,15)\) no son factibles. En efecto:
\[
2(10)+4(35)=160>120,
\]
por lo que \((10,35)\) viola la restricción de fertilizante. Además,
\[
3(30)+2(15)=120>100,
\]
por lo que \((30,15)\) viola la restricción de mano de obra.

Evaluando la función objetivo:
\[
\begin{array}{rcl}
(0,0): & 200(0)+300(0) &= 0,\\[1mm]
(0,30): & 200(0)+300(30) &= 9\,000,\\[1mm]
(20,20): & 200(20)+300(20) &= 4\,000+6\,000=10\,000,\\[1mm]
\left(\frac{100}{3},0\right): 
& 200\left(\frac{100}{3}\right)+300(0) 
&= \frac{20\,000}{3}.
\end{array}
\]
Por lo tanto, la solución óptima es
\[
(x^*,y^*)=(20,20),
\]
con ganancia óptima
\[
z^*=10\,000.
\]

Las soluciones básicas factibles en forma estándar son:
\[
\begin{array}{c|ccccc}
\text{Punto} & x & y & s_1 & s_2 & s_3 \\ \hline
1 & 0 & 0 & 45 & 100 & 120 \\[1mm]
2 & 0 & 30 & 15 & 40 & 0 \\[1mm]
3 & 20 & 20 & 5 & 0 & 0 \\[1mm]
4 & \frac{100}{3} & 0 & \frac{35}{3} & 0 & \frac{160}{3}
\end{array}
\]

\textbf{(c) Restricción Adicional: Mínimo de 25 Hectáreas de Maíz}
\vspace{2mm}
\begin{figure}
    \centering
    \includegraphics[width=0.7\linewidth]{ejercicios/Capitulo-3/ejercicio_3_3_reduccion.png}
    \caption{Reducción de la región factible}
    \label{ex:3:3:region-reducida}
\end{figure}

Si se impone la restricción adicional
\[
y\ge 25,
\]
la región factible se reduce, ya que ahora solo son admisibles los puntos que satisfacen que al menos 25 hectáreas sean plantadas con maíz.

La solución óptima anterior era
\[
(x^*,y^*)=(20,20),
\]
pero esta ya no es factible, pues
\[
20<25.
\]
Por lo tanto, el problema cambia y se debe encontrar una nueva solución óptima.

Los nuevos vértices factibles son:
\begin{itemize}
    \item \((0,25)\): intersección de \(x=0\) con \(y=25\).
    
    \item \((0,30)\): intersección de \(x=0\) con la restricción de fertilizante
    \[
    2x+4y=120.
    \]
    
    \item \((10,25)\): intersección de \(y=25\) con la restricción de fertilizante:
    \[
    2x+4y=120.
    \]
    Sustituyendo \(y=25\):
    \[
    2x+4(25)=120,
    \]
    \[
    2x+100=120,
    \]
    \[
    x=10.
    \]
    Por lo tanto, el punto es \((10,25)\).
\end{itemize}

Evaluando la función objetivo:
\[
\begin{array}{rcl}
(0,25): & 200(0)+300(25) &= 7\,500,\\[1mm]
(0,30): & 200(0)+300(30) &= 9\,000,\\[1mm]
(10,25): & 200(10)+300(25) &= 2\,000+7\,500=9\,500.
\end{array}
\]

Por lo tanto, con la restricción adicional \(y\ge 25\), la nueva solución óptima es
\[
(x^*,y^*)=(10,25),
\]
con ganancia óptima
\[
z^*=9\,500.
\]

En consecuencia, la restricción adicional modifica la solución óptima del problema. Antes se plantaban 20 hectáreas de trigo y 20 hectáreas de maíz, mientras que ahora se deben plantar 10 hectáreas de trigo y 25 hectáreas de maíz. Además, la ganancia óptima disminuye de $10\,000$
a $9\,500$
